\documentclass[11pt]{article}
\usepackage[margin=1in]{geometry}
\usepackage{amsmath,amssymb,amsthm}
\usepackage{graphicx}
\usepackage{booktabs}
\usepackage{listings}
\usepackage{xcolor}
\usepackage[numbers,sort&compress]{natbib}
\usepackage{hyperref}

\lstset{basicstyle=\ttfamily\small,breaklines=true,frame=single,columns=fullflexible,
  keywordstyle=\color{blue!60!black},commentstyle=\color{gray},language=Python}

\title{Do wildfires that start in smaller barrier compartments end smaller?\\
A re-executed pre-registered confirmatory test in northeast Washington}
\author{Justin Hart\\Viridis LLC}
\date{12 November 2026}

\begin{document}
\maketitle

\begin{abstract}
\emph{Trees Cannot Cut} holds that wildfires end where closed loops of holdable features stop them.
One consequence is that a fire which starts inside a smaller compartment of mapped barriers should end smaller.
After a pilot in the Okanogan, Methow and Chelan region of Washington State, a confirmatory replication (P1C v1.4) was
pre-registered for the adjoining northeast Washington region. Its plan, code, blind predictor table and critical
values were hash-frozen and timestamped before any fire size was requested, and the plan went through five independent audits.
The single confirmatory test regresses $\log_{10}$ final fire size (FPA FOD, 1992--2024) on $\log_{10}$ compartment
area, where compartments are bounded by TIGER roads, NHD water and LANDFIRE non-burnable fuel. It controls for the gridMET
energy-release-component percentile, fuel group, cause, month, year, barrier density, road distance and 40\,km regional fixed
effects (model M2R), and it is decided on the interval $\beta \pm 2.53\times\max(\mathrm{SE}_{\text{two-way}},\mathrm{SE}_{\text{year}})$.
We re-ran the frozen registered script, unchanged, on the registered input files, and every file's SHA-256 hash matched the registration.
The re-run gives the published numbers. On 11{,}314 ignitions, $\beta=0.0531$ (registered interval 0.0228 to 0.0835, $T=4.42$),
so the decision under the frozen rule is SUPPORTED and the pre-specified reading is REPLICATED. In plain terms, a compartment
ten times larger goes with final fires about 13\% larger.
The study's own pre-specified checks set firm limits on this result. \textbf{The study is observational. It does not show that
barriers stop fires}, and it cannot say how much of the association runs through suppression access.
\textbf{Without the 40\,km regional terms} (S13, the pilot's model M2) the estimate is 0.033, and it does not pass its critical value of 3.93.
\textbf{Lightning-caused fires show no effect} (S6a: $-0.015$, 95\% CI $-0.062$ to $0.033$); the association is in human-caused fires.
\textbf{Counting burn scars from the previous ten years as barriers} shrinks the estimate to 0.020 (S5).
\end{abstract}

\section{Introduction}

Fire managers rely on roads, water bodies and ground that will not burn as places where a fire can be held.
\emph{Trees Cannot Cut} \citep{Hart2026TCC} argues that the structure of these features matters as a network:
a fire ends where a closed loop of holdable features stops it, so the size of the loop that encloses an ignition
(its \emph{compartment}) should bound how large the fire can grow. If that holds, a fire that starts in a small
compartment should, on average, end smaller than one that starts in a large compartment, once fire weather, fuel and access are
held equal. The question matters for landscape design. If compartment size predicts fire size, the arrangement of
barriers, and not only their total amount, could be managed to limit fire spread. It also connects to operational
work on potential control locations and pre-planned containment \citep{OConnor2016,OConnor2017PCL}.

This question was put to wildfire records in October 2026 in two steps. A pilot in the Okanogan, Methow and Chelan region
(2{,}677 fires) found the association for mapped barriers. A confirmatory replication, P1C, was then designed for the
adjoining northeast Washington region. It was frozen and registered before any fire size from that region was requested
(viridis-canon pull request 138; \citealp{ViridisCanonPR138}), and was run once on 22 October 2026.

This paper reports that registered study in full and adds one element: an independent re-execution. We re-ran the frozen script,
unchanged, on the registered input files, and report whether it gives the published numbers. It does. We add no
analysis, subgroup or model. Every number below is a pre-specified output of the frozen script.

\section{Background and related work}

\paragraph{The pilot.} The pilot used FPA FOD 1992--2024 for the Okanogan, Methow and Chelan region, frozen before any
fire size was fetched. Its decisive test (H2) asked whether compartment size predicts final fire size beyond barrier
density and road distance. Under definition B (mapped barriers) it estimated $\beta=0.063$ (95\% CI 0.010 to 0.117,
$n=2{,}677$), about 16\% larger fires per tenfold larger compartment. A second definition, A50, used the Washington DNR potential control
location (PCL) raster \citep{OConnor2017PCL} thresholded at 0.5. It was underpowered and was not supported.

\paragraph{Why a confirmatory test.} The pilot's estimate cleared $p<0.05$, so it is probably inflated (winner's curse).
The P1C plan also re-read the pilot under alternative spatial error models and found its evidence only moderate. The
confirmatory test therefore needed new data, a single pre-specified test, and a critical value calibrated for spatial
correlation. The plan uses the small-telescopes benchmark of \citet{Simonsohn2015} to define what a failed replication
would mean.

\paragraph{Data sources.} Ignitions and final fire sizes come from the Fire Program Analysis fire-occurrence database
(FPA FOD), 7th edition, FPA\_FOD\_20260615 \citep{Short2026FPAFOD7}, built as described by \citet{Short2014ESSD}.
Barriers come from TIGER/Line roads, vehicle classes only \citep{CensusTIGER}; NHD large-scale water areas and waterbodies \citep{USGSNHD};
and the non-burnable classes of LANDFIRE 2024 FBFM40 \citep{LANDFIRE2024,Rollins2009LANDFIRE,ScottBurgan2005}.
Fire danger is the gridMET energy release component, expressed as a per-cell percentile over 1992--2024 \citep{Abatzoglou2013gridMET}.
Burn scars for sensitivity S5 are MTBS perimeters, used as geometry only \citep{Eidenshink2007MTBS}. Control locations for the
secondary definitions A50 and A25 come from the DNR PCL raster \citep{OConnor2017PCL}.

\section{The pre-registration}

\subsection{Freeze, registration and outcome fetch}
The plan, code, blind predictor table, raw inputs, logs and critical-value records were hash-frozen in
\texttt{FREEZE.json} and submitted to OpenTimestamps. Outcomes (fire size) were requested only after three guards passed:
G1, every frozen file unchanged; G2, an anonymous download of the public plan and \texttt{FREEZE.json} byte-identical to the local copies;
G3, no shared ignition with the pilot. The published report gives this order of events on 22 October 2026 (UTC):
registration record written at 11:30:09; addendum registration at 11:30:17; the guard download at 11:30:56; all
14{,}105 of 14{,}105 outcomes fetched by 11:34:43, none missing and none overlapping the pilot; the fetch record timestamped
at 11:34:46. \texttt{run\_test.py} then ran once. The report records no deviations from the plan beyond Addendum~1, which
changed no frozen file, model, statistic, critical value or decision rule.

\subsection{Five audits}
The plan passed through five versions, each frozen, timestamped and audited by a separate agent that had not built the study.
None was published before v1.4.
\begin{enumerate}
\item \textbf{Audit of v1.0} found an inaccurate blindness statement about MTBS attributes, a registration guard that trusted a local file,
  neglect of the 33 years shared with the pilot, unfrozen fetch evidence and undisclosed post-fire map updates.
  v1.1 fixed these and also replaced the pilot's interval with $t_{0.975,32}\times\max\mathrm{SE}$.
\item \textbf{Audit of v1.1} found that this calibration held only if spatial correlation stops at the 20\,km block edges.
  Under correlation that crosses block edges, the rule gave 6--8\% false calls. v1.2 calibrated a critical value $C_2=2.31$ under four
  spatial error models fitted to the pilot's residuals.
\item \textbf{Audit of v1.2} found that the fitted range of residual correlation (11\,km) is biased low. Under fields with ranges of 22--45\,km,
  the rule gave 6.7--9.3\% false calls. v1.3 calibrated over a screened family of error models that the pilot's residuals cannot reject (140 exponential
  fields and 9 block models), and changed the primary model to M2R, which adds 40\,km regional fixed effects. That choice was recorded before the pilot's
  M2R estimate was computed.
\item \textbf{Audit of v1.3} found that $C_2=2.37$ had no margin against anisotropic fields and compartment-level clustering, both outside the family,
  and that S13's critical value under\-covered exact long-range fields. v1.4 added a stress set of 90 structures fixed in advance, of which 52 are retained.
  This raised the critical values to $C_2=2.53$, $C_1=2.08$ and $C_{2,\mathrm{M2}}=3.93$.
\item \textbf{Audit of v1.4} (a verification audit) found no critical or moderate issue and six minor ones. Addendum~1, written and timestamped
  before any outcome was requested, answered them. It noted that between grid points a retained model could set a critical value up to
  about 0.08 higher for $C_2$ (2.61) and about 0.16 higher for $C_{2,\mathrm{M2}}$ (4.10). It therefore labels as \emph{borderline} any verdict with
  $2.53<T\le2.61$ (primary) or $3.93<T\le4.10$ (S13). These notes change no decision.
\end{enumerate}
Up to v1.2 every change widened the interval. v1.3 changed the model, and v1.4 changed only the critical values, which can only rise.
All changes were made before any outcome from the region existed.

\section{Region, data and sample}
The analysis box is longitude $-119.4$ to $-117.0$ and latitude 47.75 to 49.0. It covers eastern Okanogan, Ferry, Stevens and Pend Oreille
counties, the northern edges of Lincoln and Spokane counties, and a strip about 3\,km wide of the Idaho panhandle. The region adjoins the pilot
to the east. Predictors were computed on a 30\,m grid in EPSG:5070 over the analysis box plus 15\,km.
The sample is every FPA FOD record inside the box: 14{,}105 ignitions (4{,}779 natural, 9{,}080 human, 246 missing cause),
none shared with the pilot. After the predictor-based rules (ignitions on holdable cells and censored compartments excluded), 11{,}314 ignitions
are usable under definition B.

\paragraph{Definitions.} \textbf{B} (primary): compartments bounded by roads, water and non-burnable ground.
\textbf{A50}/\textbf{A25}: PCL $>0.5$ / $>0.25$. \textbf{B+}: B plus MTBS scars from the previous 10 years.
The PCL raster has large no-data gaps in this region, so A50 and A25 samples are small (102 and 712).

\paragraph{Known limits of blindness (disclosed in the plan).} During the pilot, MTBS perimeters with full attributes were received for
the pilot's compute box, and 63 of them reach into this region. Only geometry and year were kept. Sensitivities S10 and S10b drop the exposed ignitions.
The fuel and road maps date from 2024 but are applied to fires from 1992 to 2024. LANDFIRE updates fuels after disturbance (addressed by S11),
and post-fire roads would work against the hypothesis. Five ignitions have a pilot ignition on the same date within 5\,km; S12 drops them.

\section{Model and decision rule}
The outcome is $Y=\log_{10}(\text{fire size, acres})$. The models are
\begin{align*}
\text{M1:}\quad & Y \sim \log_{10}A_{\text{enc}} + \text{ERC}_{\text{pct}} + \text{fuel group} + \text{cause} + \text{month} + \text{year},\\
\text{M2:}\quad & \text{M1} + \log_{10}(\rho+0.001) + \log_{10}(\text{road distance}+30\,\text{m}),\\
\text{M2R:}\quad & \text{M2} + \text{40\,km region},
\end{align*}
with categorical terms as fixed effects, estimated by OLS. $A_{\text{enc}}$ is the area of the compartment that encloses the ignition and $\rho$ is barrier density.
The 40\,km regions are the cells $(\lfloor x/40000\rfloor,\lfloor y/40000\rfloor)$ in EPSG:5070, each the $2\times2$ parent of the 20\,km
clustering blocks.

The single confirmatory test is H2, definition B, model M2R: $\beta(\log_{10}A_{\text{enc}})>0$. Let $\mathrm{SE}_2$ be the CR1 standard error
two-way clustered by year and 20\,km block \citep{Cameron2011}, and $\mathrm{SE}_{\text{year}}$ the year-clustered CR1 standard error.
With $T=\beta/\max(\mathrm{SE}_2,\mathrm{SE}_{\text{year}})$, the registered interval is $\beta\pm C_2\max(\mathrm{SE}_2,\mathrm{SE}_{\text{year}})$,
$C_2=2.53$. The verdict is SUPPORTED if the interval lies above 0, CONTRADICTED if it lies below 0, and NOT SUPPORTED otherwise. The pre-specified reading is
REPLICATED if SUPPORTED. It is NOT REPLICATED if NOT SUPPORTED and the one-sided bound $\beta+C_1\max\mathrm{SE}$ ($C_1=2.08$) falls below
$\beta_{33}=0.0415$, the effect the pilot had 33\% power to detect \citep{Simonsohn2015}. It is INCONCLUSIVE otherwise.
The frozen decision code is:
\begin{lstlisting}
def decide_ci(lo, hi):
    return 'SUPPORTED' if lo > 0 else ('CONTRADICTED' if hi < 0 else 'NOT SUPPORTED')

def decide_primary(tw):
    """The registered P1C decision: the interval beta +- C2 * max(two-way SE, year-only SE)."""
    if tw.get('not_estimable'):
        return 'NOT ESTIMABLE'
    return decide_ci(tw['interval_registered'][0], tw['interval_registered'][1])
\end{lstlisting}
Across the error models in the plan's power table, the simulated probability of SUPPORTED was 0.52--0.60 at $\beta_{33}$ and 0.86--0.95 at the pilot's M2 estimate (0.063).

\section{Re-execution}
We downloaded the eleven registered input files and the frozen \texttt{run\_test.py} from the registration repository at commit
\texttt{cfdf55a24b2b833ef88dcd160813edab9c7116dd}, gunzipped the three compressed tables, and checked each SHA-256 hash against the registered list.
All twelve matched. The entry script \texttt{rerun\_p1c.py} copies the inputs into the folder layout the frozen script expects and runs it
unchanged with \texttt{python3 -B}. It then writes the full output (\texttt{p1c\_results.json}) and the primary test
(\texttt{p1c\_primary.json}). The frozen script itself refuses to run unless every input matches \texttt{FREEZE.json} and the outcomes match the
fetch record; it reported \texttt{inputs\_verified: true}.

We ran it with Python 3.12.3, numpy 2.5.4, pandas 3.1.1 and scipy 1.18.1 (the frozen script imports \texttt{scipy.stats} for a
$t$ critical value). Its only randomness, the 2{,}000-replicate year bootstrap and the 10{,}000 H4 permutations, uses a fixed generator
seed (20261020). Two runs from an empty results folder gave byte-identical files. Every number reported below agrees with the
published report to the precision it gives. The SHA-256 of our \texttt{p1c\_primary.json} differs from the hash recorded on 9 November 2026
with numpy 2.5.3 and pandas 3.0.5. To see why, we also ran the package in a scratch copy with numpy 2.5.3 and pandas 3.0.5 on our machine.
That gave a third hash. Across both runs, all decisions, readings and counts were identical, and the largest relative difference in any number
was $9\times10^{-13}$ in \texttt{p1c\_primary.json} and $6\times10^{-12}$ in \texttt{p1c\_results.json}. The hashes therefore differ only through floating-point
rounding in the last digits, which depends on the platform and linear-algebra libraries; see Section~\ref{sec:repro}.

\section{Results}

\subsection{Primary result}
\begin{table}[h]\centering
\caption{The registered primary test (H2, definition B, model M2R), from \texttt{results/p1c\_primary.json}.}
\begin{tabular}{lr}\toprule
Ignitions in the primary sample & 11{,}314 \\
Year clusters / 20\,km blocks / 40\,km regions & 33 / 75 / 24 \\
$\beta$, $\log_{10}$ size per $\log_{10}$ compartment area & \textbf{0.0531} \\
$\mathrm{SE}_2$ / $\mathrm{SE}_{\text{year}}$ / max & 0.0120 / 0.0097 / 0.0120 \\
Registered interval, $\beta\pm2.53\times\max\mathrm{SE}$ & \textbf{0.0228 to 0.0835} \\
$T=\beta/\max\mathrm{SE}$ & 4.42 \\
Decision (frozen rule) & \textbf{SUPPORTED} \\
Pre-specified reading & \textbf{REPLICATED} \\
Fire size per tenfold larger compartment & $\times1.130$ \\
\bottomrule\end{tabular}\end{table}

$T=4.42$ lies outside the borderline band $(2.53,2.61]$ of Addendum~1. Other readings are reported but do not decide.
The v1.1-rule interval ($t_{0.975,32}\times\max\mathrm{SE}$) is 0.0287 to 0.0776, which is SUPPORTED, and the pilot-rule reading is also SUPPORTED.
The year-clustered 95\% CI is 0.0342 to 0.0721 and the year cluster bootstrap gives 0.0350 to 0.0721.
For the small-telescopes reading, $\beta_{33}=0.0415$ and the 90\% upper bound is 0.078, so effects as large as those the pilot could detect are not rejected.
The barrier-density and road-distance coefficients are $-0.049$ and $0.141$.

\subsection{Pre-specified checks}
\paragraph{S13, the pilot's model M2 without regional terms ($C_{2,\mathrm{M2}}=3.93$; no decision weight).}
$\beta=0.0330$, $\max\mathrm{SE}=0.0142$, $T=2.33$. The registered interval is $-0.023$ to $0.089$, so the verdict is NOT SUPPORTED against 3.93,
and also against 4.10, so it is not borderline. Its year-clustered 95\% CI is 0.005 to 0.061 and its v1.1-rule reading is SUPPORTED.
\emph{The estimate is smaller without the 40\,km regional terms.} The plan notes that M2 is exposed to differences between reporting agencies.
It also notes that M2R compares compartments only within regions, so an effect that worked only between regions would not be detected by the primary test.

\paragraph{Leave-one-year-out.} $\beta$ ranges from 0.049 to 0.057 over the 33 years, and no year's removal changes the decision.

\paragraph{Comparison with the pilot (M2R).} The pilot's M2R estimate is 0.079, so the difference is $-0.026$ ($z=-0.98$, $p=0.33$ with two-way SEs,
from the stacked model). The pooled GLS estimate (two-way) is 0.057 (0.036 to 0.079).

\begin{table}[h]\centering\footnotesize
\caption{Sensitivity analyses, definition B. 95\% CIs; the registered-rule interval is given where the plan specifies one.}
\label{tab:sens}
\begin{tabular}{llrll}\toprule
Check & Model & $n$ & $\beta$ (95\% CI) & Registered rule \\\midrule
S10 outside the pilot's compute box & M2R & 9{,}941 & 0.043 (0.016--0.069) & SUPPORTED (0.009--0.077)\\
S10b also excl.\ pilot-fetched MTBS & M2R & 9{,}368 & 0.042 (0.012--0.072) & SUPPORTED (0.003--0.081)\\
S11 no fuel-group term & M2R & 11{,}314 & 0.051 (0.027--0.075) & SUPPORTED (0.020--0.082)\\
S12 without possible duplicates & M2R & 11{,}309 & 0.052 (0.029--0.076) & SUPPORTED (0.022--0.083)\\
S1 snap ignitions on holdable cells & M2 & 13{,}707 & 0.057 (0.039--0.075) & n/a\\
S2 include censored compartments & M2 & 11{,}585 & 0.031 (0.012--0.050) & n/a\\
S3 drop coarse coordinates & M2 & 11{,}297 & 0.033 (0.014--0.052) & n/a\\
S5 B+ (prior 10-yr burn scars) & M2 & 11{,}105 & \textbf{0.020} (0.001--0.039) & n/a\\
S6a lightning-caused only & M2 & 4{,}226 & \textbf{$-$0.015} ($-$0.062--0.033) & n/a\\
S6b human-caused only & M2 & 6{,}915 & 0.054 (0.031--0.076) & n/a\\
S8 years 2010 on & M2 & 4{,}865 & 0.049 (0.025--0.073) & n/a\\
\bottomrule\end{tabular}\end{table}

\paragraph{Secondary models (not confirmatory).}
B under M1: 0.063 (0.042--0.085), $n=11{,}314$. A50 under M1: 0.553 (0.361--0.746) and under M2: 0.527 (0.366--0.689), both $n=102$.
A25 under M1: 0.158 (0.087--0.228) and under M2: 0.140 (0.071--0.209), both $n=712$. Because of the PCL gaps the A50 sample is small. The PCL raster is a
model product that may have been trained on fire perimeters from this region.

\paragraph{H3, escaped fires} (B, M2, fires $\ge10$ acres, $n=802$): $\beta=0.091$ (0.001--0.181).
\paragraph{H4, containment} (B, $n=802$): the observed share of fires no larger than their compartment is 0.946, against a permutation null mean of 0.929
(within year $\times$ fuel-group strata), one-sided $p=0.0002$.

\section{Discussion and limitations}
Under the rule frozen and timestamped before any outcome was requested, the compartment association found in the pilot replicates in the
adjoining region. It is slightly smaller than the pilot's estimate under the same model (0.053 against 0.079), and the two estimates are statistically compatible.
In plain terms, after fire weather, fuel, cause, month, year, barrier density, road distance and 40\,km region are taken into account, fires that
start in a compartment ten times larger end about 13\% larger.

The study's own design and pre-specified checks limit what this shows:
\begin{itemize}
\item \textbf{Observational.} The study does not show that barriers stop fires. Nor can it separate fuel discontinuity from suppression access:
  roads that bound small compartments also bring crews and engines. Road distance and barrier density are controlled, but not experimentally.
\item \textbf{S13: without the regional terms the estimate is 0.033 and does not pass its critical value} (3.93). The result therefore depends on
  comparing compartments within 40\,km regions. This model was chosen, before any outcome existed, for robustness to long-range spatial correlation.
\item \textbf{S6a: lightning-caused fires show no effect} ($-0.015$, CI $-0.062$ to $0.033$). The association is in human-caused fires (S6b, 0.054).
  Human ignitions cluster near roads and settlements, where access and early detection also differ, so this pattern fits a role for
  suppression access but does not identify one.
\item \textbf{S5: counting recent burn scars as barriers shrinks the estimate to 0.020.} S2 (censored compartments included, 0.031) and S3
  (coarse coordinates dropped, 0.033) also give smaller estimates than the primary.
\item \textbf{Residual statistical risk.} The critical values cover the screened family and the stress set. They do not cover structures outside both,
  such as correlation that changes from year to year, combinations of structures at their limits, or non-exponential shapes. Between grid points the binding
  value could be up to about 0.08 higher. The primary $T=4.42$ is far from either bound.
\item \textbf{Maps and blindness.} The barrier and fuel maps are from 2024 and are applied to fires from 1992 onward. The disclosed MTBS exposure is addressed by S10 and
  S10b, which do not change the decision.
\item \textbf{Scope of the re-execution.} Our re-run shows that the frozen code gives the published numbers from the registered files. It does not check those files against the agencies'
  servers, the registration's timestamps, or the blindness of the original fetch. The frozen script checks the inputs against \texttt{FREEZE.json} and
  trusts the fetch record. The study has not been peer reviewed.
\end{itemize}

\section{Conclusions}
The pre-registered P1C confirmatory test, re-executed unchanged on the registered inputs, gives $\beta=0.0531$ (registered interval 0.0228 to 0.0835,
$n=11{,}314$). The decision is SUPPORTED and the reading is REPLICATED: fires that start in a tenfold larger compartment of mapped barriers end about 13\% larger.
This is an observational association and does not show that barriers stop fires. It weakens to 0.033, below its critical value, without the 40\,km
regional terms (S13). It is absent in lightning-caused fires (S6a), and it shrinks to 0.020 when recent burn scars count as barriers (S5).
Mechanism remains open, in particular whether suppression access rather than fuel discontinuity carries the association.

\section{Reproducibility}\label{sec:repro}
{\sloppy
The package holds \texttt{code/run\_test.py} (the frozen registered script, SHA-256 \texttt{25abae84\ldots}), \texttt{code/rerun\_p1c.py}
(the entry script, \texttt{065fe6a5\ldots}), the eleven registered inputs under \texttt{data/} with their URLs and hashes in \texttt{data\_pins.json},
and the outputs \texttt{results/p1c\_primary.json} and \texttt{results/p1c\_results.json}. From the package folder, with an empty \texttt{results/},
run \texttt{python3 code/rerun\_p1c.py}; it uses no network and takes a few minutes. It requires numpy, pandas and scipy. Our run used numpy 2.5.4, pandas 3.1.1 and scipy 1.18.1, and gave
\texttt{p1c\_primary.json} with SHA-256 \texttt{6741918a\ldots} and \texttt{p1c\_results.json} with \texttt{adda924a\ldots}.\par}
The entry script is:
\begin{lstlisting}
proc = subprocess.run([sys.executable, "-B", "run_test.py"], cwd=p1c, capture_output=True, text=True)
...
if results.get("inputs_verified") is not True or results.get("synthetic") is not False:
    sys.stderr.write("run_test.py did not verify its inputs\n")
    return 2
primary = {"registered_test": "P1C v1.4, H2, definition B, model M2R",
           "decision_rule": "SUPPORTED iff beta - C2 * max(SE_twoway, SE_year) > 0", ...}
\end{lstlisting}
The plan, addendum, audits, logs and published report are in the registration repository \citep{ViridisCanonPR138}.

\section*{Changelog}
\begin{itemize}
\item 2026-11-12 (P4-T001): first version. Re-executed the registered P1C v1.4 confirmatory test unchanged on the registered inputs and
  reported the registered study in full. No new analysis.
\end{itemize}

\bibliographystyle{plainnat}
\bibliography{refs}
\end{document}
